\RequirePackage{plautopatch}
\documentclass[platex,dvipdfmx,a4paper]{jsarticle}

\usepackage{amsmath}
\usepackage{amsfonts}
\usepackage{amssymb}

\usepackage[dvipdfmx]{graphicx}
\usepackage{float}

\usepackage{tikz}
\usepackage{caption}

\usepackage[%
  dvipdfmx,%
  setpagesize=false,%
  bookmarks=true,%
  bookmarksdepth=tocdepth,%
  bookmarksnumbered=true,%
  colorlinks=false,%
  pdftitle={},%
  pdfsubject={},%
  pdfauthor={},%
  pdfkeywords={}%
]{hyperref}
\usepackage{pxjahyper}

\usepackage{framed}
\setlength{\OuterFrameSep}{0pt}

\usepackage{physics2}
\usephysicsmodule{ab}
\usephysicsmodule{ab.braket}

\newcommand{\hb}[2]{\frac{\partial #1}{\partial #2}}
\newcommand{\bbn}[2]{\frac{d #1}{d #2}}
\newcommand{\nhb}[3]{\frac{\partial^{#1} #2}{\partial #3^{#1}}}
\newcommand{\nbbn}[3]{\frac{d^{#1} #2}{d #3^{#1}}}
\newcommand{\vectxy}[2]{\begin{bmatrix}#1\#2\end{bmatrix}}
\newcommand{\vectxyz}[3]{\begin{bmatrix}#1\#2\#3\end{bmatrix}}
\newcommand{\vect}[1]{\boldsymbol{\mathrm{#1}}}

\newcommand{\ssum}[3]{\sum_{#1=#2}^#3}

\newcommand{\veff}{V_{\mathrm{eff}}(r)}

\hypersetup{colorlinks=true}
\hypersetup{linkcolor=red}

\title{タイトル}
\author{25B00729　嶋田 悠音}

\begin{document}
\maketitle

\section*{5 双子のパラドクス}

本問では、兄が地球上に静止し、弟が一定速度 \(v\) で遠方まで移動した後、速度 \(-v\) で地球へ帰還する場合について、両者が経験する固有時を比較する。また、弟の折り返しに伴う同時性の変化をミンコフスキー図から考察する。

なお、問題文では、弟が時刻 \(t=t_0\) に地球を出発し、時刻 \(t=t_1\) に位置

\begin{align*}
  x_1=vt_1
\end{align*}

へ到達するとされている。しかし、兄の座標系で弟が時刻 \(t=t_0\) に \(x=0\) を出発し、一定速度 \(v\) で運動するならば、その世界線は

\begin{align*}
  x=v(t-t_0)
\end{align*}

である。したがって、時刻 \(t=t_1\) における弟の位置は

\begin{align*}
  x_1=v(t_1-t_0)
\end{align*}

でなければならない。問題文の \(x_1=vt_1\) は \(t_0=0\) とした場合にのみ成り立つ表式であり、一般の \(t_0\) を用いた記述とは整合しない。このため、以下では問題文の \(x_1=vt_1\) を

\begin{align*}
  x_1=v(t_1-t_0)
\end{align*}

と読み替えて議論する。

\subsection*{(1)}

まず、兄が静止している地球の慣性系を考え、横軸に空間座標 \(x\)、縦軸に時間座標 \(ct\) を取る。この座標系におけるミンコフスキー図を図\ref{fig:minkowski-coordinate}に示す。

\begin{figure}[H]
  \centering
  \begin{tikzpicture}[scale=0.9]
    \draw[->] (-4.5,0) -- (4.8,0) node[right] {$x$};
    \draw[->] (0,-0.5) -- (0,7.0) node[above] {$ct$};

    \draw[densely dashed] (0,0) -- (4.2,4.2);
    \draw[densely dashed] (0,0) -- (-4.2,4.2);

    \node[above left] at (3.7,3.7) {$x=ct$};
    \node[above right] at (-3.7,3.7) {$x=-ct$};
    \node[below left] at (0,0) {$O$};
  \end{tikzpicture}
  \caption{横軸に \(x\)、縦軸に \(ct\) を取ったミンコフスキー図}
  \label{fig:minkowski-coordinate}
\end{figure}

図中の破線は原点を通る光の世界線である。光速で \(x\) 軸の正方向または負方向へ進む光について、それぞれ

\begin{align*}
  x=ct,\qquad x=-ct
\end{align*}

が成り立つため、光の世界線は図上で傾き \(\pm 1\) の直線となる。質量を持つ物体の速度は光速を超えないため、その世界線は光の世界線よりも \(ct\) 軸に近い側に描かれる。

\subsection*{(2)}

この小問では、地球を出発して折り返し点へ向かう弟の世界線を作図する。

兄は地球上の \(x=0\) に静止しているため、その世界線は

\begin{align*}
  x=0
\end{align*}

である。一方、弟は時刻 \(t=t_0\) に地球を出発し、一定速度 \(v\) で \(x\) 軸の正方向へ進む。したがって、弟の往路の世界線、すなわち軌跡Aは

\begin{align*}
  x=v(t-t_0)
\end{align*}

で与えられる。

弟が時刻 \(t=t_1\) に折り返し点へ到達したとする。このときの位置を \(x_1\) とすれば、先に述べた読み替えにより

\begin{align*}
  x_1=v(t_1-t_0)
\end{align*}

が成り立つ。兄と弟の往路の世界線を図\ref{fig:outward-worldline}に示す。

\begin{figure}[H]
  \centering
  \begin{tikzpicture}[x=1.0cm,y=0.75cm]
    \draw[->] (-0.7,0) -- (6.2,0) node[right] {$x$};
    \draw[->] (0,-0.6) -- (0,8.8) node[above] {$ct$};

    \draw[densely dashed] (0,1.0) -- (6.0,7.0);
    \node[above left] at (5.8,6.8) {光の世界線};

    \draw[thick] (0,0) -- (0,8.2);
    \node[left] at (0,7.4) {兄の世界線};

    \draw[very thick,blue] (0,1.0) -- (4.0,6.0);
    \node[right,blue] at (2.0,3.5) {軌跡A};

    \draw[densely dashed] (0,1.0) -- (0.45,1.0);
    \draw[densely dashed] (0,6.0) -- (4.0,6.0);
    \draw[densely dashed] (4.0,0) -- (4.0,6.0);

    \fill (0,1.0) circle (2pt);
    \fill (4.0,6.0) circle (2pt);

    \node[left] at (0,1.0) {$ct_0$};
    \node[left] at (0,6.0) {$ct_1$};
    \node[below] at (4.0,0) {$x_1$};
    \node[above right] at (4.0,6.0) {$(x_1,ct_1)$};
  \end{tikzpicture}
  \caption{弟の往路における世界線}
  \label{fig:outward-worldline}
\end{figure}

図\ref{fig:outward-worldline}では、弟の世界線が光の世界線よりも \(ct\) 軸に近い側に描かれている。これは

\begin{align*}
  0<v<c
\end{align*}

に対応しており、弟の速度が光速を超えないことを表している。

\subsection*{(3)}

この小問では、弟が折り返し点から地球へ戻る世界線を作図し、兄と弟が経験する固有時を比較する。

弟は時空上の点 \((x_1,t_1)\) で進行方向を反転し、その後は一定速度 \(-v\) で地球へ戻る。したがって、弟の復路の世界線、すなわち軌跡Bは

\begin{align*}
  x-x_1=-v(t-t_1)
\end{align*}

で与えられる。

弟が地球へ帰還する時刻を \(t=t_2\) とする。帰還時には \(x=0\) であるため、軌跡Bの式に \(x=0\) と \(t=t_2\) を代入する：

\begin{align*}
  -x_1&=-v(t_2-t_1) \\
  t_2-t_1&=\frac{x_1}{v}。
\end{align*}

ここで \(x_1=v(t_1-t_0)\) を用いる：

\begin{align*}
  t_2-t_1&=t_1-t_0 \\
  t_2&=2t_1-t_0。
\end{align*}

したがって、兄の座標系では往路と復路に要する時間は等しい。兄と弟の全世界線を図\ref{fig:round-trip-worldline}に示す。

\begin{figure}[H]
  \centering
  \begin{tikzpicture}[x=1.0cm,y=0.62cm]
    \draw[->] (-0.7,0) -- (6.2,0) node[right] {$x$};
    \draw[->] (0,-0.6) -- (0,13.3) node[above] {$ct$};

    \draw[densely dashed] (0,1.0) -- (6.0,7.0);
    \node[above left] at (5.8,6.8) {光の世界線};

    \draw[thick] (0,0) -- (0,12.7);
    \node[left] at (0,9.0) {兄の世界線};

    \draw[very thick,blue] (0,1.0) -- (4.0,6.0);
    \draw[very thick,red] (4.0,6.0) -- (0,11.0);

    \node[right,blue] at (1.8,3.1) {軌跡A};
    \node[right,red] at (2.2,8.5) {軌跡B};

    \draw[densely dashed] (0,6.0) -- (4.0,6.0);
    \draw[densely dashed] (4.0,0) -- (4.0,6.0);

    \fill (0,1.0) circle (2pt);
    \fill (4.0,6.0) circle (2pt);
    \fill (0,11.0) circle (2pt);

    \node[left] at (0,1.0) {$ct_0$};
    \node[left] at (0,6.0) {$ct_1$};
    \node[left] at (0,11.0) {$ct_2$};
    \node[below] at (4.0,0) {$x_1$};

    \node[above right] at (4.0,6.0) {折り返し};
    \node[right] at (0,11.0) {再会};
  \end{tikzpicture}
  \caption{兄の世界線と弟の往復運動の世界線}
  \label{fig:round-trip-worldline}
\end{figure}

次に、兄と弟が経験する固有時を求める。固有時 \(\tau\) はミンコフスキー計量によって

\begin{align*}
  c^2d\tau^2=c^2dt^2-dx^2
\end{align*}

と定義される。弟は往路において速度 \(v\) で運動するため、

\begin{align*}
  dx=v\,dt
\end{align*}

が成り立つ。これを固有時の定義式に代入する：

\begin{align*}
  c^2d\tau^2
    &=c^2dt^2-v^2dt^2 \\
    &=c^2\left(1-\frac{v^2}{c^2}\right)dt^2。
\end{align*}

時間の正方向へ運動しているため、正の平方根を取る：

\begin{align*}
  d\tau=\sqrt{1-\frac{v^2}{c^2}}\,dt。
\end{align*}

したがって、軌跡Aに沿って弟が経験する固有時 \(\tau_A\) は

\begin{align*}
  \tau_A
    &=\int_{t_0}^{t_1}\sqrt{1-\frac{v^2}{c^2}}\,dt \\
    &=(t_1-t_0)\sqrt{1-\frac{v^2}{c^2}}。
\end{align*}

復路でも速度の大きさは \(v\) であり、兄の座標系での経過時間は \(t_2-t_1=t_1-t_0\) である。よって、軌跡Bに沿って弟が経験する固有時 \(\tau_B\) は

\begin{align*}
  \tau_B
    &=(t_2-t_1)\sqrt{1-\frac{v^2}{c^2}} \\
    &=(t_1-t_0)\sqrt{1-\frac{v^2}{c^2}}。
\end{align*}

以上より、出発から再会までに弟が経験する固有時 \(\tau_{\mathrm{弟}}\) は

\begin{align*}
  \tau_{\mathrm{弟}}
    &=\tau_A+\tau_B \\
    &=2(t_1-t_0)\sqrt{1-\frac{v^2}{c^2}}。
\end{align*}

一方、兄は \(x=0\) に静止しているため、兄については \(dx=0\) である。したがって兄の固有時は座標時間に等しく、出発から再会までに兄が経験する固有時 \(\tau_{\mathrm{兄}}\) は

\begin{align*}
  \tau_{\mathrm{兄}}
    &=t_2-t_0 \\
    &=2(t_1-t_0)。
\end{align*}

ここで \(0<v<c\) であるから、

\begin{align*}
  0<\sqrt{1-\frac{v^2}{c^2}}<1
\end{align*}

が成り立つ。よって、

\begin{align*}
  \tau_{\mathrm{弟}}<\tau_{\mathrm{兄}}
\end{align*}

を得る。したがって、再会時には弟の方が若い。

しかし、弟が往路または復路で等速運動している間だけを考えれば、弟から見て運動しているのは兄の方であり、兄の時計の方が遅れて進むように見える。この見かけ上の矛盾が双子のパラドクスである。以下では、弟が折り返す際に弟の同時刻線が変化することを用いて、この矛盾を考察する。

\subsection*{(4)}

この小問では、弟が折り返し点へ到達する直前の慣性系を考え、その慣性系において折り返し事象と同時である兄の時刻 \(t_3\) を求める。

往路における弟の慣性系を \((x',t')\) とする。兄の慣性系から弟の慣性系へのローレンツ変換は

\begin{align*}
  x'&=\gamma(x-vt), \\
  t'&=\gamma\left(t-\frac{vx}{c^2}\right)
\end{align*}

である。ただし、ローレンツ因子 \(\gamma\) は

\begin{align*}
  \gamma=\frac{1}{\sqrt{1-v^2/c^2}}
\end{align*}

である。

弟の慣性系において同時である事象は \(t'=\mathrm{const.}\) を満たす。したがって、ローレンツ変換から同時刻線の方程式は

\begin{align*}
  t-\frac{vx}{c^2}=\mathrm{const.}
\end{align*}

となる。

折り返し事象 \((x_1,t_1)\) を通る同時刻線について、定数は

\begin{align*}
  \mathrm{const.}=t_1-\frac{vx_1}{c^2}
\end{align*}

である。よって、折り返し直前の弟の同時刻線は

\begin{align*}
  t-\frac{vx}{c^2}=t_1-\frac{vx_1}{c^2}
\end{align*}

と表される。この同時刻線を図\ref{fig:outward-simultaneity}に示す。

\begin{figure}[H]
  \centering
  \begin{tikzpicture}[x=1.0cm,y=0.62cm]
    \draw[->] (-0.7,0) -- (6.2,0) node[right] {$x$};
    \draw[->] (0,-0.6) -- (0,13.3) node[above] {$ct$};

    \draw[thick] (0,0) -- (0,12.7);
    \node[left] at (0,9.2) {兄の世界線};

    \draw[very thick,blue] (0,1.0) -- (4.0,6.0);
    \draw[very thick,red] (4.0,6.0) -- (0,11.0);

    \draw[densely dashed,blue] (0,4.4) -- (5.5,6.6);
    \node[above,blue] at (2.4,5.4) {往路の同時刻線};

    \draw[densely dashed] (0,6.0) -- (4.0,6.0);
    \draw[densely dashed] (4.0,0) -- (4.0,6.0);

    \fill (0,1.0) circle (2pt);
    \fill (0,4.4) circle (2pt);
    \fill (4.0,6.0) circle (2pt);
    \fill (0,11.0) circle (2pt);

    \node[left] at (0,1.0) {$ct_0$};
    \node[left] at (0,4.4) {$ct_3$};
    \node[left] at (0,6.0) {$ct_1$};
    \node[left] at (0,11.0) {$ct_2$};
    \node[below] at (4.0,0) {$x_1$};

    \node[above right] at (4.0,6.0) {折り返し};
  \end{tikzpicture}
  \caption{折り返し直前の弟の慣性系における同時刻線}
  \label{fig:outward-simultaneity}
\end{figure}

弟が折り返す事象と同時である兄の事象は、図\ref{fig:outward-simultaneity}における同時刻線と兄の世界線 \(x=0\) との交点である。そこで、同時刻線の方程式に \(x=0\) と \(t=t_3\) を代入する：

\begin{align*}
  t_3=t_1-\frac{vx_1}{c^2}。
\end{align*}

さらに \(x_1=v(t_1-t_0)\) を代入する：

\begin{align*}
  t_3
    &=t_1-\frac{v^2}{c^2}(t_1-t_0)。
\end{align*}

したがって、弟が折り返し点へ到達する直前に、弟が折り返し事象と同時であると判断する兄の時刻は \(t_3\) である。

\subsection*{(5)}

この小問では、折り返し直後の弟の慣性系を考え、その慣性系において折り返し事象と同時である兄の時刻 \(t_4\) を求める。

折り返し直後の弟の速度は \(-v\) である。この慣性系の時間座標を \(t''\) とすれば、ローレンツ変換の時間成分は

\begin{align*}
  t''=\gamma\left(t+\frac{vx}{c^2}\right)
\end{align*}

である。したがって、折り返し直後の弟の慣性系における同時刻線は

\begin{align*}
  t+\frac{vx}{c^2}=\mathrm{const.}
\end{align*}

と表される。

折り返し事象 \((x_1,t_1)\) を通る同時刻線について、定数は

\begin{align*}
  \mathrm{const.}=t_1+\frac{vx_1}{c^2}
\end{align*}

である。よって、折り返し直後の弟の同時刻線は

\begin{align*}
  t+\frac{vx}{c^2}=t_1+\frac{vx_1}{c^2}
\end{align*}

となる。この同時刻線を図\ref{fig:return-simultaneity}に示す。

\begin{figure}[H]
  \centering
  \begin{tikzpicture}[x=1.0cm,y=0.62cm]
    \draw[->] (-0.7,0) -- (6.2,0) node[right] {$x$};
    \draw[->] (0,-0.6) -- (0,13.3) node[above] {$ct$};

    \draw[thick] (0,0) -- (0,12.7);
    \node[left] at (0,9.2) {兄の世界線};

    \draw[very thick,blue] (0,1.0) -- (4.0,6.0);
    \draw[very thick,red] (4.0,6.0) -- (0,11.0);

    \draw[densely dashed,red] (0,7.6) -- (5.5,5.4);
    \node[below,red] at (2.4,6.6) {復路の同時刻線};

    \draw[densely dashed] (0,6.0) -- (4.0,6.0);
    \draw[densely dashed] (4.0,0) -- (4.0,6.0);

    \fill (0,1.0) circle (2pt);
    \fill (4.0,6.0) circle (2pt);
    \fill (0,7.6) circle (2pt);
    \fill (0,11.0) circle (2pt);

    \node[left] at (0,1.0) {$ct_0$};
    \node[left] at (0,6.0) {$ct_1$};
    \node[left] at (0,7.6) {$ct_4$};
    \node[left] at (0,11.0) {$ct_2$};
    \node[below] at (4.0,0) {$x_1$};

    \node[above right] at (4.0,6.0) {折り返し};
  \end{tikzpicture}
  \caption{折り返し直後の弟の慣性系における同時刻線}
  \label{fig:return-simultaneity}
\end{figure}

折り返し事象と同時である兄の事象は、図\ref{fig:return-simultaneity}における同時刻線と兄の世界線 \(x=0\) との交点である。同時刻線の方程式に \(x=0\) と \(t=t_4\) を代入する：

\begin{align*}
  t_4=t_1+\frac{vx_1}{c^2}。
\end{align*}

さらに \(x_1=v(t_1-t_0)\) を用いる：

\begin{align*}
  t_4
    &=t_1+\frac{v^2}{c^2}(t_1-t_0)。
\end{align*}

したがって、折り返しの直前と直後で弟が折り返し事象と同時であると判断する兄の時刻の差は

\begin{align*}
  t_4-t_3
    &=\left(t_1+\frac{vx_1}{c^2}\right)
      -\left(t_1-\frac{vx_1}{c^2}\right) \\
    &=\frac{2vx_1}{c^2} \\
    &=\frac{2v^2}{c^2}(t_1-t_0)。
\end{align*}

折り返し前後の同時刻線を同じ図に描くと、図\ref{fig:simultaneity-jump}のようになる。

\begin{figure}[H]
  \centering
  \begin{tikzpicture}[x=1.0cm,y=0.62cm]
    \draw[->] (-0.7,0) -- (6.2,0) node[right] {$x$};
    \draw[->] (0,-0.6) -- (0,13.3) node[above] {$ct$};

    \draw[thick] (0,0) -- (0,12.7);
    \node[left] at (0,9.2) {兄の世界線};

    \draw[very thick,blue] (0,1.0) -- (4.0,6.0);
    \draw[very thick,red] (4.0,6.0) -- (0,11.0);

    \draw[densely dashed,blue] (0,4.4) -- (5.5,6.6);
    \draw[densely dashed,red] (0,7.6) -- (5.5,5.4);

    \node[above,blue] at (2.2,5.3) {折り返し直前};
    \node[below,red] at (2.2,6.7) {折り返し直後};

    \draw[<->,thick] (-0.35,4.4) -- (-0.35,7.6);
    \node[left] at (-0.35,6.0) {$c(t_4-t_3)$};

    \draw[densely dashed] (0,6.0) -- (4.0,6.0);
    \draw[densely dashed] (4.0,0) -- (4.0,6.0);

    \fill (0,1.0) circle (2pt);
    \fill (0,4.4) circle (2pt);
    \fill (4.0,6.0) circle (2pt);
    \fill (0,7.6) circle (2pt);
    \fill (0,11.0) circle (2pt);

    \node[left] at (0,1.0) {$ct_0$};
    \node[left] at (0,4.4) {$ct_3$};
    \node[right] at (0,6.0) {$ct_1$};
    \node[left] at (0,7.6) {$ct_4$};
    \node[left] at (0,11.0) {$ct_2$};
    \node[below] at (4.0,0) {$x_1$};

    \node[above right] at (4.0,6.0) {折り返し};
  \end{tikzpicture}
  \caption{弟の折り返しに伴う同時刻線の変化}
  \label{fig:simultaneity-jump}
\end{figure}

図\ref{fig:simultaneity-jump}から分かるように、弟が折り返す瞬間に、弟が兄の現在として対応させる時刻は \(t_3\) から \(t_4\) へ変化する。これは、兄の時計が実際に瞬間的に進むことを意味しない。弟が往路の慣性系から復路の慣性系へ移ることによって、空間的に離れた事象に対する同時性の定義が変化することを意味している。

したがって、「弟から見れば兄の時計の方が遅れる」という議論を往路と復路について単純に接続することはできない。弟は折り返しによって慣性系を変更しており、その際に兄の世界線上で同時と判断する事象が変化する。よって、兄と弟の立場は旅行全体を通して対称ではなく、再会時に弟の方が若いという結論と矛盾しない。以上より、通常の双子のパラドクスは同時性の相対性を考慮することで解消される。

\subsection*{(6)}

この小問では、空間が \(x\) 方向に長さ \(L\) の周期的境界条件を持つ場合を考える。すなわち、ある慣性系において

\begin{align*}
  (x,t)\sim(x+L,t)
\end{align*}

とし、同じ時刻 \(t\) における位置 \(x\) と位置 \(x+L\) を同一の時空上の点として扱う。

この場合、弟は折り返すことなく一定速度 \(v\) で進み続けることにより、空間を一周して兄と再会できる。兄の慣性系において、弟が空間を一周するために必要な時間を \(T\) とする。弟が進む距離は \(L\) であるため、

\begin{align*}
  vT=L
\end{align*}

が成り立つ。したがって、

\begin{align*}
  T=\frac{L}{v}
\end{align*}

を得る。この状況を、周期的空間を平面上に展開した図\ref{fig:periodic-worldline}に示す。

\begin{figure}[H]
  \centering
  \begin{tikzpicture}[x=0.9cm,y=0.75cm]
    \draw[->] (-0.7,0) -- (7.2,0) node[right] {$x$};
    \draw[->] (0,-0.6) -- (0,8.3) node[above] {$ct$};

    \draw[thick] (0,0) -- (0,7.8);
    \draw[thick] (6.0,0) -- (6.0,7.8);

    \draw[very thick,blue] (0,0.8) -- (6.0,6.8);

    \draw[densely dashed] (0,6.8) -- (6.0,6.8);
    \draw[densely dashed] (6.0,0) -- (6.0,6.8);

    \fill (0,0.8) circle (2pt);
    \fill (6.0,6.8) circle (2pt);

    \node[below] at (0,0) {$0$};
    \node[below] at (6.0,0) {$L$};
    \node[left] at (0,5.2) {兄の世界線};
    \node[above left,blue] at (4.0,4.8) {弟の世界線};
    \node[right] at (6.0,6.8) {再会};
    \node[left] at (0,6.8) {$cT$};

    \draw[<->] (0,-0.35) -- (6.0,-0.35);
    \node[below] at (3.0,-0.35) {$L$};

    \draw[->,thick] (6.35,2.3) .. controls (7.05,2.9) and (7.05,4.7) .. (6.35,5.3);
    \node[right,align=left] at (6.85,3.8) {$x=L$を\\$x=0$と同一視};
  \end{tikzpicture}
  \caption{周期的境界条件を持つ空間における兄と弟の世界線}
  \label{fig:periodic-worldline}
\end{figure}

兄は静止しているため、再会までに兄が経験する固有時は

\begin{align*}
  \tau_{\mathrm{兄}}=T=\frac{L}{v}
\end{align*}

である。一方、弟は一定速度 \(v\) で運動しているため、弟が経験する固有時は

\begin{align*}
  \tau_{\mathrm{弟}}
    &=T\sqrt{1-\frac{v^2}{c^2}} \\
    &=\frac{L}{v}\sqrt{1-\frac{v^2}{c^2}}。
\end{align*}

よって、

\begin{align*}
  \tau_{\mathrm{弟}}<\tau_{\mathrm{兄}}
\end{align*}

が成り立ち、再会時には弟の方が若い。

この場合、兄と弟はともに途中で加速度運動を行っていないため、局所的な運動状態だけを比較すれば、両者は対等な慣性観測者であるように見える。しかし、仮定した周期的境界条件は、すべての慣性系において同じ形を取る条件ではない。

このことをローレンツ変換によって確かめる。兄の慣性系において同一視される二つの事象を

\begin{align*}
  (x,t),\qquad (x+L,t)
\end{align*}

とする。弟の慣性系 \((x',t')\) へのローレンツ変換は

\begin{align*}
  x'&=\gamma(x-vt), \\
  t'&=\gamma\left(t-\frac{vx}{c^2}\right)
\end{align*}

である。ただし、

\begin{align*}
  \gamma=\frac{1}{\sqrt{1-v^2/c^2}}
\end{align*}

である。

まず、二つの事象の弟の慣性系における空間座標の差を求める：

\begin{align*}
  \Delta x'
    &=\gamma\bigl((x+L)-vt\bigr)-\gamma(x-vt) \\
    &=\gamma L。
\end{align*}

次に、時間座標の差を求める：

\begin{align*}
  \Delta t'
    &=\gamma\left(t-\frac{v(x+L)}{c^2}\right)
      -\gamma\left(t-\frac{vx}{c^2}\right) \\
    &=-\gamma\frac{vL}{c^2}。
\end{align*}

したがって、兄の慣性系における同一視

\begin{align*}
  (x,t)\sim(x+L,t)
\end{align*}

は、弟の慣性系では

\begin{align*}
  (x',t')
  \sim
  \left(
    x'+\gamma L,\,
    t'-\gamma\frac{vL}{c^2}
  \right)
\end{align*}

と表される。

すなわち、兄の慣性系では同時刻にある二つの空間点が同一視されているが、弟の慣性系では、同一視される二つの事象の時刻が

\begin{align*}
  \gamma\frac{vL}{c^2}
\end{align*}

だけ異なる。このため、周期的境界条件は弟の慣性系では、兄の慣性系と同じく空間座標だけの周期条件として表すことができない。

特殊相対性原理は、すべての慣性系が物理的に対等であり、物理法則がすべての慣性系で同じ形を取ることを要請する。しかし、本問で仮定した

\begin{align*}
  (x,t)\sim(x+L,t)
\end{align*}

という空間方向だけの周期的境界条件は、特定の慣性系でのみ時間のずれを伴わない形を取る。したがって、この境界条件は、その条件を単純な形で記述できる兄の慣性系を特別な慣性系として選んでいる。

よって、空間方向だけに周期的境界条件を課した時点で、時空全体に対する特殊相対性原理は成り立たなくなる。ただし、各点の近傍ではミンコフスキー計量が保たれており、局所的な物理法則に対するローレンツ不変性まで失われるわけではない。

以上より、周期的境界条件を持つ場合に兄と弟の年齢差が生じても矛盾はないが、これは兄と弟が特殊相対性原理の下で対等であるにもかかわらず年齢差が生じることを意味するのではない。空間方向だけの周期的境界条件が大域的に特別な慣性系を選び、特殊相対性原理の大域的な対称性を破っているため、兄と弟の立場が非対称になっているのである。
\end{document}